% -----------------------------------------------------------------------------------------------------
% B3. Statements restated verbatim from theory/B3_prop59_tightness.tex (Appendix B3), in the order of the note.
% Numbers quoted (each with its setting):
%   ratio tau*/tau_book of the note's paragraph "Comparison with the textbook choice" (Remark rem:B3-conclusion):
%     in [1.05, 3), at most 4/3 for kappa >= 9n/16; for n = 1000: 32/(3(5+sqrt(17))) = 1.169 at kappa = 1e3,
%     2.42 at kappa = 1e2, 2.94 at kappa = 10; speed-ups 3.2--3.7 for kappa <= 1e2 (digest.md, B3 tables).
%     The joint optimum is at a smaller step (gamma* R^2 = 0.081 at kappa = 1e2, 0.0098 at kappa = 10); at the textbook
%     step the best Delta = c gamma^2 (c = 5.04, 5.89) gains 1.81 and 1.98 (note, Outcome; Section sec:exp-constants).
%   B3_summary.json, problems: theta_0 = 0; ||eta*|| = 1.0--1.3 at kappa = 10 and >= 5.9 for kappa >= 1e2 (R_x = 1);
%     lmin(F''(eta*))/lambda from 1.0003 to 69.9, below 2 in 9 of the 15 configurations (the six with kappa <= 1e2,
%     and kappa = 1e3 with alpha = 1, 2 and kappa = 1e4 with alpha = 2, where the region rate gains 1.37, 1.01, 1.07).
% -----------------------------------------------------------------------------------------------------
\subsection{Proposition~5.9 as a chain of inequalities}\label{sec:th-B3}

Hypothesis B3 asks which of the inequalities I1--I6 of the proof of Proposition~\ref{prop:59} explain the gap between
its certified contraction factor $\rho_{\rm bound}=1-\min\{\frac1{3n},\frac{3\mu}{16R^2}\}$ and the contraction observed
along SAGA. This subsection states the instruments of the diagnosis (an exact one-step identity, a tightness ratio for
each inequality, and an attribution of the gap to the inequalities) and delimits what a better analysis could certify:
the best constants of the proof, a lower bound on the exponent of $\kappa$, and a trajectory-dependent version of the
rate.

\paragraph{Exact identity and tightness ratios.}
Fix $t\ge1$ and write $(\theta,z)=(\theta_{t-1},z^{(t-1)})$ for the state before the $t$-th draw and
$\mathcal F_{t-1}$ for the $\sigma$-field generated by $i(1),\dots,i(t-1)$, so that $\Ei=\E[\,\cdot\mid\mathcal F_{t-1}]$.
Let $D_{t-1}=\norm{\theta-\etastar}^2$, $M_{t-1}=\sum_i\norm{z_i-\fip(\etastar)}^2$,
$G_{t-1}=\inner{\Fp(\theta)}{\theta-\etastar}$, $S_{t-1}=\sum_i\norm{\fip(\theta)-\fip(\etastar)}^2$,
$\Omega_{t-1}=\Ei\norm{\omega_t}^2$ and $B_{t-1}=\Ei\norm{b_t}^2=\frac1n\sum_i\norm{\fip(\etastar)-z_i+\bar z}^2$, where
$\bar z$ is the mean of the $z_i$ and $b_t$ is as in I1; the index $t-1$ is dropped when no confusion is possible. With a
free weight $\Delta>0$ in place of $3\gamma^2$ in \eqref{eq:lyapunov}, $V_{t-1}=D_{t-1}+\Delta M_{t-1}$. All these
quantities are explicit averages over the $n$ values of $i(t)$. Before any inequality is applied, the conditional
expectation of the Lyapunov function is an exact function of them, which serves as the reference for every certified
factor:
\restate{prop:B3-exact}
Thus $\rho_{\rm exact}(t)$ is computable from the state at the cost of one pass over the $n$ components, and
Section~\ref{sec:exp-B3} records it along SAGA trajectories. In this notation I1--I4 read $\Omega\le\frac2nS+2B$,
$\frac1nS\le R^2G$, $B\le\frac1nM$ and $G\ge\mu D$, and I5 replaces the two coefficients left by I1--I4 by the larger
one. Each inequality receives a ratio, its smaller side over its larger side, equal to $1$ exactly when the inequality
is an equality at the current state:
\restate{def:B3-ratios}
The effective constants $\muEff=G/D$ and $\LEff=S/(nG)$ turn I4 and I2 into identities at the current iterate. By
\resref{lem:B3-ranges}, for $\theta\ne\etastar$ the ratios of I1, I2, I4 and I5 lie in $(0,1]$, $r_{\rm var}$ lies in
$[0,1]$ and $\mu\le\muEff\le\LEff\le R^2$ ($r_{\rm init}\in(0,1]$ when $\theta_0\ne\etastar$); moreover $r_{\rm sq}=\frac12$
exactly when the gradient changes $a_i=\fip(\theta)-\fip(\etastar)$ and the memory errors $b_i=\fip(\etastar)-z_i+\bar z$
are uncorrelated under $i(t)$. As \resref{rem:B3-fresh} explains, one expects $r_{\rm sq}$ close to $\frac12$ when the
memory is stale and $r_{\rm sq}\le\frac12$, possibly $\ll\frac12$, when it is nearly fresh.

\paragraph{Representation and attribution.}
Inserting the ratios into \eqref{eq:B3-EV} gives the two coefficients of $\Ei[V_t]$ exactly, and shows what the proof
does: it substitutes worst-case values for four quantities that depend on the state.
\restate{prop:B3-representation}
With the textbook constants, the parameter rate $1-c_D$ is $\frac{\muEff}{2R^2}$ times a bracket in $[\frac38,1)$, so
that $r_{\rm coco}$ and $r_{\rm sq}$ change it by at most a factor $\frac83$, and the memory rate $1-c_M$ lies between
$\frac1{3n}$ and $\frac1n$, so that $r_{\rm sq}$ and $r_{\rm var}$ change it by at most a factor $3$; only $\muEff/\mu$,
which lies in $[1,\kappa]$ by \resref{lem:B3-ranges} and is therefore not bounded independently of $\kappa$, and the
step I5, which replaces the weighted mean $(c_DD+c_M\Delta M)/V_{t-1}$ by $\max\{c_D,c_M\}$,
can produce larger factors. Undoing the substitutions one at a time, in either of two orders, gives monotone chains from the
certified to the exact factor, hence an attribution of the gap to the individual inequalities:
\restate{prop:B3-attribution}
The increments $\nu_k$ add up to the gap between the exact and the certified contraction exponents, and the
\emph{share} of an inequality is its increment divided by this gap. The shares depend on the order. Order~B gives each
inequality its full effect on the weighted contraction; in order~A, undoing I4 and I2, which only lower $c_D$, has no
effect as long as $c_M$ is the larger coefficient ($\kappa<\frac{9n}{16}$ for the textbook constants), and undoing I3,
which only lowers $c_M$, has none as long as $c_D$ is. By \resref{rem:B3-orders}, comparing the two orders shows which
component of $V_t$ is the bottleneck of the bound. Section~\ref{sec:exp-B3} measures the ratios and the shares along
SAGA trajectories on logistic regression.

\paragraph{Optimal constants and the exponent of $\kappa$.}
The identities and I1--I3 do not involve $\gamma$ or $\Delta$, so the same chain certifies a rate for every pair
$(\gamma,\Delta)$ of an explicit admissible set; optimizing over this set measures how much of the gap is due to the
constants $\gamma=\frac1{4R^2}$, $\Delta=3\gamma^2$ of Proposition~\ref{prop:59}.
\restate{prop:B3-general}
Writing $\Delta=c\gamma^2$ separates the two constants: the memory rate depends on $c$ only, and for fixed $c$ the
parameter rate is a concave function of the step.
\restate{cor:B3-c}
Since $c>2$, no weight is admissible when $\gamma R^2\ge\frac12$. At the textbook step every $c\in(2,6)$ is admissible,
and the best of these weights certifies a rate asymptotically equivalent to $\frac2{3n}$, twice the textbook
$\frac1{3n}$, as $n/\kappa\to\infty$ (Section~\ref{sec:B3-constants}). Maximizing $\tau(c)$ over $c$ gives the best
constants of the proof structure in closed form:
\restate{prop:B3-optimum}
These are the optimal rate and step of \cite[Theorem~1 and Corollary~1]{hofmann2015variance} for SAGA: their Lyapunov
function has the same two-term form, but its memory term averages upper bounds $H_i\ge\norm{z_i-\fip(\etastar)}^2$
obtained from Bregman divergences, and their proof uses the same I1 and I3 but applies strong convexity and
cocoercivity through function values; the two chains differ, but their optima coincide (Section~\ref{sec:B3-constants}).
The ratio of $\tau^\star$ to the textbook rate $\min\{\frac1{3n},\frac3{16\kappa}\}$ lies in $[1.05,3)$. For
$\kappa\ge\frac{9n}{16}$ it is at most $\frac43$: there the textbook constants are within $\frac43$ of the best of the
proof structure (the ratio is $\frac{32}{3(5+\sqrt{17})}\approx1.17$ for $n=1000$, $\kappa=10^3$;
Section~\ref{sec:exp-constants}). For $n=1000$ and $\kappa\le10^2$ the constants matter more: the optimal constants, at a
smaller step, raise the certified rate by a factor $2.42$ ($\kappa=10^2$) or $2.94$ ($\kappa=10$), of the same size as the
speed-ups $3.2$--$3.7$ of the exact over the certified exponent measured there at the textbook step
(Section~\ref{sec:exp-B3}), as noted in \resref{rem:B3-conclusion}. At the textbook step itself, the best weight $\Delta$
gains a factor $1.81$ or $1.98$ (Section~\ref{sec:exp-constants}). This optimality is relative to the proof structure
(Lyapunov function $D+\Delta M$, inequalities I1--I5): another Lyapunov function or another chain could certify more.
Whatever the analysis, however, this Lyapunov function cannot certify more than $\frac1n$ per iteration: the memory
recursion \eqref{eq:B3-EM} of \resref{lem:B3-identities}, $\Ei[M_t]=(1-\frac1n)M_{t-1}+\frac1nS_{t-1}$, gives
$\Ei[M_t]\ge(1-\frac1n)M_{t-1}$, hence $\E[V_t]\ge\Delta(1-\frac1n)^tM_0$, and no bound $\E[V_t]\le c_0\rho^tV_0$ valid for
all $t$, with $c_0$ independent of $t$ and $M_0>0$, has $1-\rho>\frac1n$, whatever $\Delta>0$. When
$\kappa\le\frac{9n}{16}$, where the textbook rate is $\frac1{3n}$, a factor $3$ is thus the most that any analysis of
$\E[V_t]$ can gain.

Conversely, the exponent $1$ of $\kappa$ is intrinsic to SAGA with a constant step, not to its analysis. The following
instance satisfies \resref{ass:B3-setting}, the setting of Section~\ref{sec:setting}, and \eqref{eq:B3-saga} in its
statement is the SAGA iteration~\eqref{eq:saga}:
\restate{prop:B3-lower}
The argument only uses that the direction is unbiased and that $\Fp$ is affine: on any quadratic $F$ the mean SAGA
iterate is the gradient-descent iterate with the same step, whatever the memory, and the variance only adds to the
error. For $\kappa\gg n$ the textbook bound is within the factor $\frac{64}3$ of this bound on its leading term, and within
$\frac83$ of the bound $(2\kappa-\frac12)\log\frac1{2\varepsilon}$ that the same argument gives for the textbook step;
caveat~(ii) of \resref*{rem:B3-conclusion} derives both comparisons.

\paragraph{Trajectory-dependent version.}
Strong convexity I4 and cocoercivity I2 are only used at the current iterate, so the proof also certifies rates in which
effective constants replace $\mu$ and $R^2$; this is the hook for hypothesis B2 (Section~\ref{sec:th-B2}). Regard
$\muEff(\theta)=\inner{\Fp(\theta)}{\theta-\etastar}/\norm{\theta-\etastar}^2$ and
$\LEff(\theta)=\frac1n\sum_i\norm{\fip(\theta)-\fip(\etastar)}^2/\inner{\Fp(\theta)}{\theta-\etastar}$ as functions of
$\theta\ne\etastar$ (they do not depend on the memory), and write $\muEff(t-1)=\muEff(\theta_{t-1})$ and
$\LEff(t-1)=\LEff(\theta_{t-1})$:
\restate{prop:B3-traj}
In the second bound, which uses both effective constants, $\muEff$ enters linearly, whereas a small $\LEff$ helps by at
most a factor $\frac83$, as explained in \resref{rem:B3-fixed-step}. These factors are random ($\mathcal F_{t-1}$-measurable) and correlated with
$V_{t-1}$, so the conditional bounds cannot be iterated: the product of the factors along a trajectory, which
Figure~\ref{fig:fig1} plots, does not bound $\E[V_T]$ in general. An a-priori rate needs a deterministic bound on the
effective constants over a region $\Theta$ of the parameter space, and follows from the supermartingale
$\rho_\Theta^{-(t\wedge\tau)}V_{t\wedge\tau}$, where $\tau$ now denotes the exit time from $\Theta$ (not a rate as in
$\tau(c)$ and $\tau^\star$):
\restate{prop:B3-region}
For logistic regression, the self-concordance-type bound $\abs{\ell'''}\le\ell''$ \cite{bach2010self,bach2014adaptivity}
gives an explicit lower bound on $\mu_\Theta$ over balls around $\etastar$:
\restate{prop:B3-logistic}
Applied from the first time the iterate enters the ball of radius $\frac3{2R_x}$ around $\etastar$ (by the strong Markov
property of $(\theta_t,z^{(t)})$, with the Lyapunov value at that time in place of $V_0$), \resref{prop:B3-region}
gives, until the next exit, a contraction factor at most
$1-\min\{\frac1{3n},\frac{3}{32R^2}(\lambda+\lmin(F''(\etastar)))\}$: the parameter term of Proposition~\ref{prop:59} is
multiplied by $\frac12(1+\lmin(F''(\etastar))/\lambda)$, which changes the rate only when
that term binds ($\kappa>\frac{9n}{16}$). From $\theta_0=0$, the starting point lies outside this ball in all our
problems with $\kappa\ge10^2$ ($\norm{\etastar}\ge5.9$, $R_x=1$). In the fifteen configurations of
Section~\ref{sec:exp-B3}, $\lmin(F''(\etastar))/\lambda$ ranges from 1.0003 to 69.9 and is below 2 in nine of them: six
have $\kappa\le10^2$, where the memory term binds and the rate is unchanged, and the other three gain less than a
factor $\frac32$.

\paragraph{Known and new.}
Proposition~\ref{prop:59} and its proof are those of \cite{bach2024learning}; Lyapunov functions adding memory errors, or upper
bounds on them, to the squared distance are used in \cite{hofmann2015variance,hu2017unified}; the optimal constants of
this proof structure (\resref*{prop:B3-optimum}) are those of \cite{hofmann2015variance}; for oblivious step schedules
depending only on $R^2$, the exponent $1$ of $\kappa$ follows from \cite[Theorem~2]{arjevani2016iteration}. To our
knowledge, the use of the exact identity as an instrument, the reading of the proof as a substitution of worst-case
values with monotone attribution chains, the elementary lower bound for every $n$ and every constant step, and the
trajectory-dependent bounds are not in the sources we reviewed in this form.
